\section{Results and Discussion}
\label{sec:results}

This section presents qualitative and quantitative analysis of the Joint-Embedding Predictive and Flow Matching architectures across 1,300 benchmark evaluations spanning 11 physical and cyber-physical domains.

\subsection{Performance Across IoT and Remote Sensing Domains}
As shown in Table~\ref{tab:multidataset_benchmark}, the flow-matching joint-embedding models achieve top performance across all four IoT and remote sensing domains:

\begin{enumerate}[leftmargin=*]
    \item \textit{Domain 1: Satellite and Planetary Remote Sensing (SMAP, MSL, SWAN -- 78 channels):} On NASA Earth observation satellite and Mars rover telemetry, \texttt{dyflow\_jepa} and \texttt{ncad\_dyflow\_jepa} achieve the highest Point-F1 across all evaluated models ($0.285$ and $0.294$ on SMAP; $0.265$ and $0.254$ on MSL), outperforming TranAD ($0.183$ and $0.206$) and TimesNet ($0.178$ and $0.174$). On SMAP channel \texttt{A-7}, \texttt{ncad\_dyflow\_jepa} achieves $0.9841$ Point-F1 with $1.0000$ precision ($0$ false alarms) and $0.9688$ recall, demonstrating that Optimal Transport velocity modeling tracks orbital telemetry mode shifts. On SWAN coastal remote sensing, \texttt{ncad\_dyflow\_jepa} reaches $0.629$ Point-F1, compared to TranAD ($0.354$) and TimesNet ($0.261$).
    \item \textit{Domain 2: Wearable IoT and Biomechanical Edge Sensing (Daphnet -- 26 channels):} On wearable tri-axial accelerometry nodes tracking Parkinsonian gait freezing, \texttt{patch\_jepa} ($0.275$ Point-F1, $0.557$ PA-F1) and \texttt{patch\_flow} ($0.235$ Point-F1, $0.532$ PA-F1) outperform all published baselines (TimesNet: $0.146$, DCdetector: $0.118$, TranAD: $0.103$). On channel \texttt{S08R01E3}, \texttt{patch\_jepa} reaches $0.7624$ Point-F1 with $90.26\%$ recall, while \texttt{patch\_flow} achieves a $1.0000$ Oracle ceiling on \texttt{S08R01E1}.
    \item \textit{Domain 3: Industrial IoT and Cyber-Physical Systems (GHL, Genesis, GECCO -- 16 channels):} In multi-phase gas-oil loop telemetry, subtle temperature and level drifts cause standard contrastive and reconstructive models to yield low scores (\texttt{dcdetector}: $0.000$, \texttt{ncad}: $0.005$). In contrast, \texttt{patch\_flow} ($0.225$ Point-F1) and \texttt{patch\_jepa} ($0.190$ Point-F1) maintain high sensitivity and achieve $100\%$ recall across multiple fault channels with low false alarm rates.
    \item \textit{Domain 4: Smart Infrastructure and Edge Environmental Networks (10 channels):} On network flow telemetry subject to DDoS bursts and port scans (\texttt{cicids}), \texttt{ncad\_dyflow\_jepa} achieves $0.260$ Point-F1. On Channel 7, it achieves $0.8988$ Point-F1 with $99.03\%$ precision ($89,179$ true positives vs $878$ false alarms) and $0.9960$ PA-F1, whereas TimesNet fires $6,413$ false alarms and DCdetector scores $0.000$. On smart environmental sensing (\texttt{room-occupancy}), \texttt{patch\_flow} achieves $0.368$ Point-F1. On the highway traffic stream (\texttt{metro}), all tested architectures achieve near-zero Point-F1 ($0.000-0.004$) due to high entropy and low signal-to-noise ratio in sparse traffic volume.
\end{enumerate}

\subsection{Comparison with Baselines}
As detailed in Table~\ref{tab:multidataset_benchmark}, DyFlow-JEPA achieves the highest aggregate scores across both aggregation schemes:
\begin{itemize}[leftmargin=*]
    \item \textit{Channel-Weighted Micro Evaluation (130 Channels):} \texttt{flow\_jepa} achieves the highest score of $0.2389$ Point-F1, representing a $+52.5\%$ improvement over TimesNet ($0.1566$), $+56.2\%$ over TranAD ($0.1529$), $+146\%$ over DCdetector ($0.0971$), and $+216\%$ over Anomaly Transformer ($0.0756$).
    \item \textit{Unweighted Dataset Macro Evaluation (11 Domains):} Flow matching models (\texttt{flow\_jepa}: $0.1921$, \texttt{ncad\_flow\_jepa}: $0.1922$) consistently lead the macro evaluation, outperforming TranAD ($0.1610$), TimesNet ($0.1510$), and Anomaly Transformer ($0.1083$).
\end{itemize}

Key comparative findings include:
\begin{enumerate}[leftmargin=*]
    \item \textit{DyFlow-JEPA vs. DCdetector (Contrastive Learning):} While DCdetector \cite{yang2023dcdetector} performs moderately on periodic signals, it experiences scores of $0.0000$ Point-F1 on non-stationary, sparse channels across GHL, Genesis, and cicids. In contrast, DyFlow-JEPA's non-contrastive manifold regularization avoids dimensional collapse across all 11 domains.
    \item \textit{DyFlow-JEPA vs. Anomaly Transformer \& TimesNet (Reconstructive Models):} Reconstructive architectures \cite{xu2022anomaly, wu2023timesnet} optimize raw signal reconstruction $x \to \hat{x}$. On non-stationary channels, they trigger frequent false alarms (e.g., TimesNet produces $76,264$ false positives on \texttt{cicids-0} and $55,186$ on \texttt{GHL-17}). By operating in representation space, DyFlow-JEPA filters high-frequency sensor noise and models physical dynamics directly.
    \item \textit{Additive Value of Velocity Field Modeling:} The progression from base \texttt{ncad} ($0.0859$) $\to$ \texttt{point\_jepa} ($0.2137$) $\to$ \texttt{flow\_jepa} ($0.2389$) confirms that continuous Optimal Transport velocity modeling provides substantial gains over deterministic regression.
\end{enumerate}

\subsection{Point Adjustment vs. Point-Level Evaluation}
A key finding from our 1,300-run evaluation is the empirical disparity between legacy Point-Adjusted F1 (PA-F1) and strict Point-Level F1 (Point-F1). Under PA-F1, baselines like TimesNet and TranAD report inflated scores ($0.2853$ and $0.2517$) because single incidental spikes within an anomaly window turn entire segments into true positives. However, under unadjusted Point-F1, their precision drops severely due to tens of thousands of false positives. 

On datasets with isolated single-timestep transition events or abrupt binary regime switches (such as \texttt{room-occupancy} and \texttt{Genesis}), unadjusted Point-F1 and PA-F1 mathematically coincide (e.g., $0.338/0.338$ for DyFlow-JEPA on \texttt{room-occupancy} and $0.144/0.144$ for Patch-Flow on \texttt{Genesis}) because the ground-truth anomaly events contain no multi-step segments to be expanded by point adjustment. While unadjusted Point-F1 scores in the $0.20-0.35$ range reflect the difficulty and low base-rate of exact timestep anomaly detection in multi-sensor telemetry, DyFlow-JEPA improves relative performance across all tested domains.

\subsection{Phase-Space Geometric Attractor Dynamics}
The behavior of the joint-embedding representations is visualized in the 2D and 3D phase-space attractor trajectories of Figure~\ref{fig:phase_space_attractor} for Point-JEPA and Figure~\ref{fig:dyflow_jepa_validation}(a) for DyFlow-JEPA's continuous tangent velocity vector field. In nominal operation, physical dynamics map onto a closed, smooth limit-cycle manifold, where the DyFlow-JEPA velocity field $v_\psi(z_t, t, z_{\text{ctx}})$ anticipates trajectory evolution with near-zero discrepancy $S(t) \approx 0$. When physical faults occur, the state diverges out of the periodic attractor into a degenerate trajectory, generating a large directional velocity discrepancy vector that breaches the EVT decision threshold with minimal latency.
